 \subsection{Constructing cuspidal ribbons}\label{cuspcons}
For \(\beta \in \Psi\), we write \(\init(\beta) \subseteq \N\) for the set of nodes \(b \in \N\) such that \(\beta = \alpha(\res(b), \height(\beta))\). Note then that \(\init(\delta) = \N\), and if \(\beta \in \Phi_+^\re\) then \(\init(\beta) = \N_t\) for some \(t \in \Z_e\).
 

\begin{defi}\label{defxib}
Let \(\beta \in \Psi\) and \(b \in \init(\beta)\). Define a \(\no/\ea\) path \((z^i)_{i=1}^{\height(\beta)}\) by setting \(z^1 = b\), and 
\begin{align*}
z^i = \begin{cases}
\no z^{i-1} & \textup{if } \alpha(\res(b), i-1) \succ \beta;\\
\ea z^{i-1} & \textup{if } \beta \succ \alpha(\res(b), i-1),\\
\end{cases}
\end{align*}
for \(i = 2, \ldots, \height(\beta)\). Then define \(\zeta^{(\beta, b)} = \{z^i\}_{i=1}^{\height(\beta)}\).
\end{defi}

\begin{lemma}\label{zetaiscusp}
The set of nodes  
\(\zeta^{(\beta,b)}\) is a cuspidal ribbon of content \(\beta\), with \(\zeta^{(\beta,b)}_{\ssww} = b\).
\end{lemma}
\begin{proof}
First, note that the path \((z^i)_{i=1}^{\height(\beta)}\) is well-defined, which follows from Lemma~\ref{notapprox}\eqref{cond:notapprox2}. We have that \(\zeta^{(\beta,b)}\) is a ribbon by Lemma~\ref{hookispath}, and \(\cont(\zeta^{(\beta,b)}) = \beta\) by Lemma~\ref{hookcont}.

It remains to show that \( \zeta = \zeta^{(\beta,b)}\) is cuspidal. Let \(\nu \subsetneq \zeta\) be an \(\ssee\)-removable ribbon in \(\zeta\). By Lemma~\ref{CuspForSkewHook}, it will be enough to show that \(\cont(\nu) \succ \beta\). By Lemma~\ref{xiuv}, we have 
\(\nu = \xi^\zeta_{z^i,z^j}\) for some \(1 \leq i \leq j \leq h\), where \(\so z^i, \ea z^j \notin \zeta\). 


First assume \(i=1, j<h\). Since \(\ea z^j  \notin \zeta\), it follows from the construction of \(\zeta\) that \(z^{j+1} = \no z^j\), and thus \(\alpha(\res(b),j) \succ \beta\). But then by Lemma~\ref{hookcont} we have 
\begin{align*}
\cont(\nu) = \cont(\xi^\zeta_{z^1, z^j}) =  \alpha(\res(b),j) \succ \beta,
\end{align*}
as desired.

Now assume \(1<i\). Since \(\so z^i \notin \zeta\), we have \(z^i = \ea z^{i-1} \) and thus \(\beta \succ \alpha(\res(b), i-1)\). If \(j<h\), then since \(\ea z^j  \notin \zeta\), it follows from the construction of \(\zeta\) that \(z^{j+1} = \no z^j\), and thus \(\alpha(\res(b),j) \succ \beta\). On the other hand, if \(j=h\), then \(\alpha(\res(b),j) = \beta\), so in any case we have \(\alpha(\res(b),j) \succeq \beta\).
Therefore, applying Lemma~\ref{hookcont}, we have
\begin{align*}
\alpha(\res(b), j) &= 
\cont(\xi^\zeta_{z^1, z^j}) 
=\cont(\xi^\zeta_{z^1, z^{i-1}}\sqcup \xi^\zeta_{z^i, z^j})\\
&=
\cont(\xi^\zeta_{z^1, z^{i-1}}) + \cont(\xi^\zeta_{z^i, z^j})
= \alpha(\res(b), i-1) + \cont(\nu).
\end{align*}
If \(\beta \succeq \cont(\nu)\), we would have by Lemma~\ref{GenConvexLem}\eqref{cond:genconvex3} that \(\beta \succ \alpha(\res(b), j)\), a contradiction. Thus \(\cont(\nu) \succ \beta\), as desired.
\end{proof}


\begin{lemma}\label{onlycuspribbon}
Let \(\nu\) be a ribbon of content \(\beta \in \Psi\). Then \(\nu\) is cuspidal if and only if \(\nu = \zeta^{(\beta, \nu_{\ssww})}\).
\end{lemma}
\begin{proof}
The `if' direction is proved in Lemma~\ref{zetaiscusp}. Now we prove the `only if' direction. Let \(b =\nu_{\ssww}\), and \(h = \height(\beta)\). By Lemma~\ref{hookcont} we must have \(\beta = \alpha(\res(b),h)\). Consider the skew shape \(\zeta = \zeta^{(\beta,b)}\), and let \((z^i)_{i=1}^h\) be the \(\no/\ea\) path constructed in Definition~\ref{defxib} so that \(\zeta = \{z^i\}_{i=1}^h\). By Lemma~\ref{hookispath}, there is a \(\no/\ea\) path \((w^i)_{i=1}^h\), with \(w^1 = b\) and \(\nu = \{w^i\}_{i=1}^h\). Assume by way of contradiction that \(\nu \neq \zeta\). Then there exists \(2 \leq t \leq h\) such that \(w^i = z^i\) for \(i<t\), and \(w^t \neq z^t\). Note that by Lemma~\ref{notapprox}, we have \(\beta \not \approx \alpha(\res(b),t-1)\).

First assume \(\beta \succ \alpha(\res(b),t-1)\). Then \(z^t = \ea z^{t-1}\).
We have \(w^t \in \{\no w^{t-1}, \ea w^{t-1}\} = \{\no z^{t-1}, \ea z^{t-1}\}\), but \(w^t \neq z^t\), so \(w^t = \no z^{t-1}\). Then \(\ea w^{t-1} \notin \nu\), and \(\ea w^h \notin \nu\), so \(\xi^\nu_{w^1, w^{t-1}}\) is \(\ssee\)-removable in \(\nu\) by Lemma~\ref{xiuv}. By cuspidality of \(\nu\) and Lemma~\ref{CuspForSkewHook} we have that
\begin{align*}
\alpha(\res(\nu_1), t-1) \succ \cont(\xi^\nu_{\nu_t, \nu_h}) \succ \beta,
\end{align*}
a contradiction.

Now assume \(\alpha(\res(b),t-1) \succ \beta\). Then \(z^t = \no z^{t-1}\), and we have \(w^t \in \{\no w^{t-1}, \ea w^{t-1}\} = \{\no z^{t-1}, \ea z^{t-1}\}\), but \(w^t \neq z^t\), so \(w^t = \ea z^{t-1}\). Then \(\so w^t \notin \nu\), and \(\ea w^h \notin \nu\), so \(\xi^\nu_{w^t, w^h}\) is \(\ssee\)-removable in \(\nu\) by Lemma~\ref{xiuv}. By cuspidality of \(\nu\) and Lemma~\ref{CuspForSkewHook} we have that
\begin{align*}
\alpha(\res(w^t), h-t + 1) = \cont(\xi^\nu_{w^t, w^h}) \succ \beta.
\end{align*}
But then by convexity and (\ref{splithookcont}) we have
\begin{align*}
\beta = \alpha(\res(b), h) =  \alpha(\res(b),t-1) + \alpha(\res(w^t), h-t + 1) \succ \beta,
\end{align*}
a contradiction. Thus, in any case we derive a contradiction, so we must have \(\nu = \zeta\) as desired.
\end{proof}


\subsection{Distinguished cuspidal ribbons}
For each \(\beta \in \Phi_+^\re\), we make a distinguished choice of \(b^\beta \in \init(\beta)\). For each \(t \in \Z_e\), we make a distinguished choice of \(b^t \in \N_t\). Then we define the distinguished cuspidal ribbons, for all \(\beta \in \Phi_+^\re\) and \(t \in \N_t\):
\begin{align*}
\zeta^\beta := \zeta^{(\beta, b^{\beta})}, \qquad \zeta^t := \zeta^{(\delta, b^t)}
\end{align*}

By construction, the path associated with \(\zeta^{(\beta, b)}\) in Definition~\ref{defxib} relies only on \(\res(b)\) and \(\beta\), and not on the specific location of \(b\). Thus by Lemmas~\ref{zetaiscusp} and~\ref{onlycuspribbon} we have the following result:

\begin{prop}\label{allcusprib}
For all \(\beta \in \Phi_+^\re\), \(t \in \ZZ_e\), we have
\begin{align*}
[\zeta^{\beta}] = \{\zeta^{(\beta,b)} \mid b \in \init(\beta)\},
\qquad
\textup{and}
\qquad
[\zeta^{t}] = \{\zeta^{(\delta,b)} \mid b \in \N_t\}.
\end{align*}
The set \(\{\zeta^{\beta} \mid \beta \in \Phi_+^\re\} \cup \{\zeta^t \mid t \in \Z_e\}\) represents a complete and irredundant set of cuspidal ribbons with content in \(\Psi\), up to \(e\)-similarity.
\end{prop}
